\subsection*{把效果上限分成两段检查}

从空购买库与平台余额 10000 出发，先看本地受理。完整 \(S\) 唯一，P1、C1 共同提交，一笔 P1 只有一个 C1。受理后的原输入和支付输入固定，所有发送读取原映射。两次新确认分别产生 \(S_1\)、\(S_2\) 和 C1、C2，平台身份也分别不同。

再看平台。设 \(P(S)\) 是 \(S\) 关联的本地购买集合，\(E(C)\) 是固定 C 对应的成功收款数。当每个可能生效的到达都受平台去重保护，原映射持续保留时，一次全额支付模型希望保持
\[
 |P(S)|\leq 1,\qquad E(C)\leq 1.
\]
前一项依赖本地身份仲裁。要得到后一项，平台对固定范围与键提供的效果上限，还要与本地 C 的固定映射接起来。只有本地唯一性，P1 可以只有一行，X1、X2 却各扣一次。平台若只保证每个键至多成功一次，本地给 C1 换两个键，两次调用仍会被当成独立操作，因而不能推出 \(E(C)\leq 1\)。所有发送读取同一映射，才把这两段仲裁连起来。

\CSevenFigure{ch07-counts}{本地购买数与外部收款数由不同的仲裁限制}{fig:ch07-counts}

平台资金还满足自身账本关系。沿用余额记号，以 \(b_0\) 表示平台付款账户初始余额，\(b\) 表示当前余额，\(D\) 为已提交的平台扣款流水集合，\(a(d)\) 是一条流水的正整数金额。只允许这些购买支出时，
\[
 b=b_0-\sum_{d\in D}a(d),\qquad b\geq 0.
\]
正确重试留下 C1、X1 和一条 3000 分流水，代入为 \(10000-3000=7000\)。换键的错误历史有两条各 3000 的流水，结果为 4000。后者可以保持资金守恒，却违反同 C1 的效果上限，所以核对余额和总额以后，还要沿身份关系数出效果。

本地 \code{paid} 也有对应事实：证据关联 C1、X1 和约定金额，同时存在归属唯一的 O1 与完整保存结果；等待付款或最终拒绝没有 O1。反方向容许间隔：X1 已成功，本地还没有 R。这正是开头的未决状态，恢复路径以后补齐结论。若要求两库状态始终相等，就会把实际允许的故障历史排除在模型外。

逐步检查这些关系：A 只增加受理和固定映射；第一次 E 在平台事务中共同增加收款、流水、原结果并减少余额；同键重复到达找到原记录，保留原效果；R 根据核对后的证据发布本地结论，重复、旧观察和矛盾终态各走上一节的分支。应用在步骤之间退出，改变的是后续还缺哪份证据，已提交记录由所属系统继续保留。

\subsection*{按下游能力选择继续方式}

表~\ref{tab:ch07-capabilities} 比较相同故障下的继续路径。表内幂等受理均以固定身份、正确输入和有效保护窗口为条件；查询以恢复者有定位信息、了解读取范围为条件。

\begin{table}[htbp]
\centering\small
\caption{不同下游能力对应的继续路径}\label{tab:ch07-capabilities}
\begin{tabularx}{\linewidth}{@{}p{3.1cm}XX@{}}
\toprule
下游能力 & 未决 C1 怎样继续 & 能形成什么结论\\
\midrule
幂等受理及可信结果查询 & 所有生效到达受保护时同身份发送，由查询或可信通知补结论 & 平台效果可限制至多一次，证据可得时确定结局\\
幂等受理，重放可解释终态 & 保护期间同参重发，读取原终态 & 限制效果并恢复结论，依赖完整返回机会\\
幂等受理，只重放未决错误 & 同参发送保持上限，等待其他证据 & 效果上限可成立，最终获知依赖其他通道\\
只有可信查询 & 查已有结果，空结果保持未决 & 找到记录可确认，延迟旧请求仍可能执行\\
两项都缺少 & 保留原映射及未决记录，进入可核查的人工或账单流程 & 自动继续缺少同时保证安全与进展的依据\\
\bottomrule
\end{tabularx}
\end{table}

可信查询可以先报告处理中，随后再报告已完成，关键是每个状态与时点都有可解释规则。搜索存在延迟时，找到并核对成功可以提供正面证据，空结果的证明力较弱。仅查询的平台增加原子关闭或可证明的请求失效能力以后，就改变了表内的前提，可以重新推导安全重发。

本地 \code{paid} 表示已经保存的证据，平台成功表示外部业务事实，两者之间可以有间隔。记录保留、路由正确、平台结果可访问，恢复任务也得到运行时，查询、回调或对账可以使本地补齐。若把这段等待直接写成失败，就会在 E 与 R 之间发布错误终态。

\clearpage
\subsection*{自动进展还需要持续运行的恢复者}

有效幂等键让重复发送安全，结束业务还依赖执行机会。服务在 A 后永久停止，平台从未收到请求，效果上限为零也满足安全条件；E 已发生但结果消息永久丢失，本地又可能一直等待。两种停滞位于不同阶段。

对尚未执行的 C1，一组足够条件包括：原输入和映射持续可读；恢复者持续取得运行机会；发送仍在平台保护范围内；网络与平台最终允许某次请求被受理；平台处理能够到达可解释结论。最终成功还依赖付款授权、资金与平台业务条件，明确业务拒绝也可以是正确结束。

对已经执行的 C1，本地最终保存结论还依赖至少一个可用结果通道、仍可定位的原身份、保留到被读取的证据，以及能够提交 R 的本地数据库。客户端获知又需要一次完整结果到达并保存。外部执行、本地记录和客户端获知，各自有完成条件。

\CSevenFigure{ch07-progress-evidence}{E 形成外部收款，R 保存本地结论，客户端获知依赖完整结果到达并保存}{fig:ch07-progress-evidence}

固定保护窗口进一步约束这些条件。任意长的网络中断可能跨过安全发送期限，本地继续保留受理事实，之后沿查询或账单核查。AWS Builders' Library 讨论迟到请求时，把知识保留与资源生命周期及额外窗口联系起来，并单独讨论同一标识携带不同参数~\cite{aws-idempotent}。实际迟到范围、平台规则和证据保留共同决定窗口，本地延长重试天数不会改变平台能力。

P1 的受理、C1 的固定映射、允许发送的条件以及终态证据已经有了具体规则。进程在 A 后退出，接下来要让另一执行者重新找到 C1，判断何时查询或发送。原执行者恢复后，也要知道自己还能修改哪些记录。下一章从已经保存的 P1、C1 出发，建立持久恢复任务与接管规则。
